\section{The cyclotomic quartic as a worked example}
\label{sec:cyclotomic}

For $f(T)=T^4+1$, the preceding argument gives an explicit density.
The arithmetic of $\Q(\zeta_8)$ also makes several parts of the
determinant construction particularly concrete. We describe that
specialization and two additional ways of organizing its estimates:
the level equation on integral representatives and the negative Pell
equation at the upper end of the prime range.
We conclude with a conditional estimate for analytic relations.
The factorization viewpoint is closely related to Heath-Brown's
Gaussian treatment of $n^2+1$ and Reuss's number-field construction
\cite{HeathBrown2012Squarefree,Reuss2015PowerFree}.

\subsection{Local factors and the dyadic conclusion}

\begin{corollary}\label{cor:cyclotomic-density}
Put
\[
 c_8=\prod_{p\equiv1\pmod8}\left(1-\frac4{p^2}\right).
\]
Then $c_8>3/4$, and
\[
 \#\{1\le n\le X:n^4+1\text{ is squarefree}\}=c_8X+o(X).
\]
Consequently, for every sufficiently large real $N$,
\[
 \#\{n\in\Z:N\le n\le2N,\ n^4+1\text{ is squarefree}\}
 \ge N/2.
\]
\end{corollary}

\begin{proof}
The polynomial $(T+1)^4+1$ is Eisenstein at $2$, so $T^4+1$ is
irreducible. Its value at zero is one, which proves local
admissibility. There is no root modulo $4$: the value is $1$ for
even inputs and $2$ for odd inputs. For an odd prime $p$, a root of
$T^4+1$ has order eight in $\mathbb F_p^\times$. Conversely, the cyclic
group $\mathbb F_p^\times$ has precisely four elements of order eight
when $p\equiv1\pmod8$, and none otherwise. Each root is simple, since
$4T^3$ is nonzero there, and therefore lifts uniquely modulo $p^2$.
Theorem~\ref{thm:density} gives the displayed Euler product and asymptotic.

Every prime in this product is at least $17$. For finite products,
$\prod_j(1-a_j)\ge1-\sum_j a_j$ when $0\le a_j\le1$; passing to
the limit gives
\[
 c_8\ge1-4\sum_{p\equiv1\pmod8}p^{-2}
 >1-4\sum_{m\ge17}m^{-2}
 >1-4\int_{16}^{\infty}t^{-2}\,dt=\frac34.
\]
Taking the difference of the asymptotics at $2N$ and $N$, with an
$O(1)$ endpoint correction, proves the final assertion.
\end{proof}

The central prime range admits a stronger error estimate than the
positive-proportion conclusion by itself requires. We retain it
explicitly, including the small extension below the input scale.

\begin{proposition}\label{prop:cyclotomic-central}
With $\epsilon_0=10^{-6}$, there exists $\delta_0>0$ such that
\[
 \#\{n\in\Z:N\le n\le2N,\ p^2\mid n^4+1\text{ for some prime }
       N^{1-\epsilon_0}\le p\le N^{3/2+\epsilon_0}\}
 \ll N^{1-\delta_0}.
\]
\end{proposition}

\begin{proof}
Proposition~\ref{prop:large-prime-tail} already treats the primes
$p>N$. For the remaining range, split into dyadic intervals with
lower endpoint $P=N^\eta$, where $1-\epsilon_0\le\eta\le1$.
Use the first quartic parameter cell with the following enlarged
endpoint bounds:
\[
 \eta_-=1-\epsilon_0,\quad \eta_+=\frac{3001}{3000},\quad
 b_-=\frac{2999}{1500},\quad b_+=2+2\epsilon_0,
 \quad t_* =\frac12,\quad w=\frac{24809}{100000}.
\]
Keep mode $\mathsf Q$ and $\tau=1/50000$. These are fixed rational
parameters. In the notation of Section~\ref{sec:parameters}, they give
\[
 m=\frac{74427}{100000},\quad
 u=\frac{76819}{1200000},\quad
 v=\frac{313933}{1000000}.
\]
Both strict cut inequalities \eqref{eq:parameter-cut-inequalities}
hold, $m<\min(\eta_-,b_-)$, and $3/20\le v/u\le5$.
An exact check is supplied in Appendix~\ref{app:quartic-extension}.
The conservative upper bounds for the thresholds in decreasing
dimension order are
\[
 (\overline\delta_4,\overline\delta_3,
   \overline\delta_2,\delta_1)
 =\left(\frac{4729}{100000},\frac{5097}{100000},
         \frac{6899}{100000},\frac{76819}{1200000}\right).
\]
Taking cumulative maxima gives
\[
 1-m-\sum_{h=1}^4\overline E_h-4\tau
   =\frac{1941}{100000}>0,\qquad
 1-m-2u-2\tau=\frac{15319}{120000}>0.
\]
Thus the lattice summation and infinite-fiber deletion have fixed
positive savings on this entire enlarged cell.

We check the arithmetic hypotheses rather than inferring them from
the numerical test. Since $\epsilon_0<1/2$, every prime now considered
satisfies $p^2\ge N^{2-2\epsilon_0}>N$. It also tends to infinity,
so the fixed bad primes are eventually absent. These are the only
uses of $p>X$ in the proof of
Lemma~\ref{lem:balanced-factorization}: exclusion of the fixed bad
primes and at most one input per root class modulo $p^2$ in an
interval of length $N$. Repeating that proof therefore gives the
same factorization and selected-point multiplicities in this band.
The paired charts, determinant cuts and lattice groups require only
the displayed coordinate upper bounds and their strict parameter
inequalities. The Hilbert and curve alternatives concern these
arithmetic equations and weights and impose no further condition
$p>N$. The argument of Section~\ref{sec:tail-proof} in the
number-field range consequently gives one power saving on every
dyadic interval here. Delete the one global quintic-value set before
this summation, absorb the $O(\log N)$ intervals by a smaller saving,
and combine with the already proved $p>N$ bound. The endpoint
correction at $n=N$ is $O(1)$.
\end{proof}

\subsection{An explicit factorization and a Pell estimate}

Write $\zeta=\exp(2\pi i/8)$ and $R=\Z[\zeta]$. In this case the
ideal-class correction in Lemma~\ref{lem:balanced-factorization} can
be omitted: one can take $\mu=1$ and obtain primitive coefficient
vectors. The following direct argument records why.

The ring $R=\Z[i]+\zeta\Z[i]$ is Euclidean for the absolute field
norm. Given $a+b\zeta\in\Q(\zeta)$, round $a,b\in\Q(i)$ to
Gaussian integers, with errors $c,d$ satisfying
$|c|^2,|d|^2\le1/2$. Then
\[
 |N_{\Q(\zeta)/\Q}(c+d\zeta)|=|c^2-id^2|^2<1.
\]
The triangle inequality gives strictness unless both errors have
squared modulus $1/2$. In that remaining case their real and
imaginary parts have absolute value $1/2$, so $c^2$ and $id^2$ are
perpendicular and the squared modulus is $1/2$. Euclidean division
implies that every ideal of $R$ is principal.

If $p^2\mid n^4+1$, evaluation $\zeta\mapsto n$ defines
surjections $\phi_j:R\to\Z/p^j\Z$ for $j=1,2$. Write
$\ker\phi_1=(A)$. Its index is $p$, so $N(A)=p$; the norm is
positive because the embeddings occur in complex conjugate pairs.
The inclusion $(A)^2\subseteq\ker\phi_2$ is an equality by their
indices. It follows that
\begin{equation}\label{eq:cyclotomic-factorization}
 A^2B=n-\zeta,\qquad A,B\in R,\qquad N(A)=p.
\end{equation}
The unit $1+\sqrt2$, where $\sqrt2=\zeta-\zeta^3$, has absolute
values $1+\sqrt2$ and $(1+\sqrt2)^{-1}$ at the two complex places.
Multiplication by a suitable power of this unit balances the two
absolute values of $A$. Thus, if $N\le n\le2N$ and $P\le p<2P$,
\[
 |\sigma(A)|\asymp P^{1/4},\qquad
 |\sigma(B)|\asymp NP^{-1/2}
\]
at every embedding. Both coefficient vectors in the basis
$1,\zeta,\zeta^2,\zeta^3$ are primitive, since a common integer
divisor of either vector would divide the coefficient $-1$ of
$\zeta$ in $A^2B$.

\begin{lemma}\label{lem:cyclotomic-pell}
Uniformly in the positive integer $D$, the number of positive integer
pairs $(n,q)$ with $N\le n\le2N$ and $n^4+1=Dq^2$ is
$O(1+\log N)$. Consequently, for every fixed $\epsilon>0$, the
number of inputs in this interval with $p^2\mid n^4+1$ for some
prime $p>N^{3/2+\epsilon}$ is
$O(N^{1-2\epsilon}(1+\log N))$.
\end{lemma}

\begin{proof}
Put $x=n^2$. If $D$ is a square, factoring
$Dq^2-x^2=1$ rules out $x>0$. Otherwise attach to a solution the
positive number $\alpha=x+q\sqrt D=x+\sqrt{x^2+1}$, of norm $-1$.
For two distinct solutions with $\alpha_2>\alpha_1$, their ratio
$\gamma=\alpha_2/\alpha_1$ belongs to $\Z[\sqrt D]$ and has norm
one, because $\alpha_1^{-1}=-\overline{\alpha_1}$. Writing
$\gamma=a+b\sqrt D$ gives
$2a=\gamma+\gamma^{-1}>2$, whence $a\ge2$ and
$\gamma\ge2+\sqrt3$. All the $\alpha$ are at most $8N^2+1$.
Their fixed multiplicative spacing proves the first assertion.
For the second, $D=(n^4+1)/p^2\le17N^{1-2\epsilon}$; sum the first
bound over these integers $D$.
\end{proof}

A fixed projective $B$-direction has at most two primitive integral
representatives and hence fixes $D=N(B)$. Taking norms in
\eqref{eq:cyclotomic-factorization} then gives the same Pell equation
with $q=p$. A fixed $A$-direction fixes $p=N(A)$ and one residue
class for $n$ modulo $p^2$, so it contributes
$O(1+N/p^2)$ inputs. These are special projective-direction
multiplicity bounds. The general affine argument instead selects one
representation per input and uses the fixed-$x$ bound in
Lemma~\ref{lem:balanced-factorization}.

\subsection{Directions and the level equation}

For coefficient vectors $A,B$, let $T_0,S,C_2,C_3$ be the coefficients
of $1,\zeta,\zeta^2,\zeta^3$ in $A^2B$. Their projective directions
lie on
\[
 \mathcal V_8=\{C_2=C_3=0\}\subseteq\mathbb P^3\times\mathbb P^3,
\]
while the actual integral representatives satisfy
\begin{equation}\label{eq:cyclotomic-level}
 T_0=n,\qquad S=-1.
\end{equation}
If $H,J$ denote the hyperplane classes of the two factors, the
complete-intersection calculation in Lemma~\ref{lem:bicone} becomes
\[
 [\mathcal V_8]=(2H+J)^2,\qquad
 (H^4,H^3J,H^2J^2,HJ^3,J^4)_{\mathcal V_8}=(0,1,4,4,0).
\]
On the split-coordinate torus, with conjugates $\zeta_i$, its points
satisfy $A_i^2B_i=T_0+S\zeta_i$. Thus the divisor $S=0$ is, on that
open set, the inverse-square graph $[B_i]=[A_i^{-2}]$.
The analytic coordinate $-S/T_0$ is exactly $1/n$ at the counted
points. This identifies the concrete direction geometry behind the
paired charts of Section~\ref{sec:arithmetic-boxes}.

There is also a determinant saving obtained directly from $S=-1$.
We give its column-selection form; it explains why that equation
contains more information than its projective zero locus.

Let $D\subseteq\mathcal V_8$ be an integral direction variety of
dimension $r\ge1$ meeting $\{S\ne0\}$, and write
$\mathcal R_D(K,L)$ for the restrictions of
bihomogeneous forms of bidegree $(K,L)$. Suppose a homogeneous
polynomial $v_0(k,l)$ of degree $r$, with nonnegative mixed-degree
coefficients, supplies the lower bound
\begin{equation}\label{eq:cyclotomic-volume-input}
 \dim\mathcal R_D(Tk,Tl)
 \ge\frac{v_0(k,l)}{r!}T^r-O(T^{r-1})
\end{equation}
on each fixed positive rational ray. One can use the actual
mixed-degree polynomial, or coordinatewise lower bounds for its
coefficients. For bounded geometric degree, the error can be taken
uniformly on any fixed finite list of rays. Indeed, after the
corresponding projective embedding, linear Noether normalization of
the coordinate ring gives a polynomial subring whose field extension
degree is the projective degree $e$. Monomials of degree at most
$e-1$ contain a field basis: their spans increase strictly until they
fill the extension. After multiplication by powers of one normalization
coordinate, these give a graded injection of $e$ copies of that
polynomial ring, shifted by $e-1$. Its degree-$T$ dimension is
$e\binom{T-e+1+r}{r}$, which proves the required uniform lower bound.

Suppose the integral representatives in polynomial-height lattice
bases satisfy $|A|\ll N^h$, $|B|\ll N^g$, where $h,g\ge0$.
For positive rational $k,l$, put $a=\min(k/2,l)$ and assume
$v_0(k,l)>0$. A nested selection of integral columns has average
logarithmic size, divided by the common bidegree scale $T$, at most
\begin{equation}\label{eq:cyclotomic-level-saving}
 kh+lg-(2h+g)\int_0^a
       \frac{v_0(k-2j,l-j)}{v_0(k,l)}\,dj+\varepsilon.
\end{equation}
Here $\varepsilon>0$ is arbitrarily prescribed: first fix a sufficiently
fine mesh and then a sufficiently large $T$, before $N$ grows.

To prove this, fix a rational mesh $0=j_0<\cdots<j_s<a$ and consider
the nested subspaces
\[
 S^{j_iT}\mathcal R_D(T(k-2j_i),T(l-j_i))
 \subseteq\mathcal R_D(Tk,Tl).
\]
Multiplication by $S$ is injective, since $S$ is nonzero on $D$.
Choose decreasing integers
$r_i=v_0(k-2j_i,l-j_i)T^r/r!+O(T^{r-1})$ below their dimensions,
and put $P=r_0$. Starting from the smallest subspace, extend to $r_i$ independent
generators at each stage. Each generator is an individual integral
form $S^b\mathcal M$ of bidegree $(Tk,Tl)$. Its value at a counted
point is $(-1)^b\mathcal M(A,B)$, bounded by
\[
 O_T\bigl(N^{h(Tk-2b)+g(Tl-b)}\bigr).
\]
At least $r_i$ columns have $b\ge j_iT$. Summing these exponents,
then refining the fixed mesh, gives
\eqref{eq:cyclotomic-level-saving}. This argument does not require
small coefficients for the rebased form $S$.

For the local divisibility step, suppose that $D$ is defined over $\Q$
and has an integral model smooth of relative dimension $r$ at the
chosen residue class modulo $q$. A determinant of $P$ such
common-bidegree columns in that projective residue class has
valuation at least $s_r(P)$, as in
\eqref{eq:affine-local-orders}: normalize by chart coordinates that
are units modulo $q$ and expand in the $r$ free differences.
Combining that divisibility with
\eqref{eq:cyclotomic-level-saving}, a prime $q\ge N^\rho$ forces a
proper cut whenever
\begin{equation}\label{eq:cyclotomic-level-threshold}
 \rho>\frac{1+1/r}{v_0(k,l)^{1/r}}
 \left\{kh+lg-(2h+g)\int_0^a
       \frac{v_0(k-2j,l-j)}{v_0(k,l)}\,dj\right\}.
\end{equation}
The strict margin fixes the mesh and determinant degree before $N$.
In particular, for a curve with $a_1=H\cdot D$, $b_1=J\cdot D$,
the choice $(k,l)=(2,1)$ makes the integral $1/2$ and the threshold
$(2h+g)/(2a_1+b_1)$. This is the projective counterpart of the
weighted affine curve estimates in Section~\ref{sec:curves}; the
general proof uses those affine estimates throughout.

\subsection{Analytic relations and box incidence}

The main proof uses an exact polynomial relation in each pure direction
component. We record a local analytic version with its hypotheses,
describe the quintic equality case, and give an elementary box-incidence
estimate. These observations explain uses of the
cyclotomic geometry; the proof of
the density theorem uses the affine constructions of
Sections~\ref{sec:determinant-cuts} and~\ref{sec:lattice-boxes}.

\paragraph{Approximate rank from an actual analytic relation.}
Fix $\omega>0$ and give four local variables $u_1,u_2,u_3,u_4$ weights
$(\omega,\omega,\omega,1)$. For the paired quartic charts one takes
$\omega=w$, the grid exponent, whereas $m=3w$ counts the boxes.
Let $f_\lambda$ be holomorphic on one fixed complex polydisc about
zero, uniformly bounded there, with Taylor coefficients continuous at
a parameter $\lambda_0$. Assume that $f_{\lambda_0}$ is a nonzero germ
of lowest weighted degree $D>0$. Consider points satisfying
\[
 |u_i|\ll N^{-\omega}\ (1\le i\le3),\qquad
 |u_4|\ll N^{-1},\qquad f_\lambda(u)=0,
\]
where $\lambda$ lies in a sufficiently small neighborhood of
$\lambda_0$. For each evaluation matrix, fix one such $\lambda$;
all its points satisfy the same equation $f_\lambda=0$.
All truncation levels and matrix sizes below are fixed
before $N$ grows. If the limiting germ instead has nonzero constant
term, there are no such shrinking zeros near $\lambda_0$ for large $N$.

Let $E_Z$ be the coefficient space spanned by monomials of weight
strictly less than $Z$, with its Euclidean coefficient norm. For fixed
$Z>D$, truncated multiplication defines
\[
 M_\lambda:E_{Z-D}\longrightarrow E_Z,\qquad
 p\longmapsto[p f_\lambda]_{<Z}.
\]
This map is injective at $\lambda_0$. Indeed the lowest weighted
parts of two nonzero polynomials have nonzero product, and the lowest
weight of $p f_{\lambda_0}$ is less than $Z$ for nonzero
$p\in E_{Z-D}$. A nonzero maximal minor persists near $\lambda_0$,
so $M_\lambda$ has a uniformly bounded left inverse on its image.

At the chosen points, every truncated multiple
$[u^\alpha f_\lambda]_{<Z}$ is $O_Z(N^{-Z})$. To justify this uniformly,
expand on the common polydisc to an ordinary Taylor order $K$ with
$\min(\omega,1)(K+1)\ge Z$. The omitted finite terms have weight at
least $Z$, and the remaining Taylor tail has the same bound. The
identity $f_\lambda(u)=0$ then gives the assertion for the truncated
multiple. The bounded left inverse gives that bound for every
unit-norm element of $W_\lambda=\operatorname{im}M_\lambda$.

Now take any fixed number of analytic column functions, uniformly
bounded on the same polydisc, and evaluate them at the chosen points.
Truncate each column to $E_Z$ and project its coefficient vector onto
$W_\lambda^\perp$. The evaluation matrix changes entrywise by
$O_Z(N^{-Z})$, while the projected matrix has rank at most
\[
 q_Z=\dim E_Z-\dim E_{Z-D}
     =\frac{D}{6\omega^3}Z^3+O_{D,\omega}(Z^2+1).
\]
For the last equality choose a lowest monomial $u^\gamma$ of the
limiting relation, so its weight is $D$. The dimension difference
counts the monomials of weight less than $Z$ not divisible by
$u^\gamma$. Their exponents lie in the union of the strips
$\alpha_i<\gamma_i$. The individual strips have total leading
coefficient $D/(6\omega^3)$, and pairwise intersections contribute
$O(Z^2+1)$.

This is a bound on an approximating matrix, not an exact rank bound
on the analytic evaluations. For an $\ell$-by-$\ell$ evaluation matrix
$G$, the singular values in decreasing order satisfy
$\sigma_{q_Z+1}(G)\ll_{\ell,Z}N^{-Z}$ when $q_Z<\ell$.
Choosing finitely many levels $Z$ and multiplying these bounds gives
\[
 |\det G|\ll_\ell
 N^{-\frac34(6\omega^3/D)^{1/3}\ell^{4/3}+O(\ell)}.
\]
Indeed the $j$th singular value may use
$Z=(6\omega^3j/D)^{1/3}-O(1)$, and summing these levels gives the
displayed exponent. The finitely many small indices for which this
level is at most $D$ use the uniform $O(1)$ bound on the matrix entries
and are absorbed in $O(\ell)$. The neighborhoods are chosen for those finitely
many levels. Thus no rate of convergence of $\lambda$ to $\lambda_0$
is needed. A compact parameter family of nonzero pulled-back germs,
with the same analytic bounds and coefficient continuity, gives uniform
constants by a finite cover; if their orders vary, use the largest
order among the finitely many centers. Normalizing the original
sections alone does not establish this nonvanishing condition on
the pulled-back germs.

\paragraph{The quintic equality case.}
The polynomial line curves of Section~\ref{sec:curves} also have a
concrete interpretation here.
If their first coordinates are linear in a rational parameter and
their second coordinates polynomial, then
\[
 (t-\tau_i)^2\mid n(t)-\zeta_i\qquad(1\le i\le4).
\]
The least possible degree of $n(t)$ is five. After normalizing two
critical points to $0,1$, Proposition~\ref{prop:quintic-discard}
places these quintics in one fixed finite family. Here the splitting
field is $\Q(\zeta)$ itself, of degree four. Integrality of a quintic
value places its input in a fixed fractional lattice, with all
embedding sizes $O(N^{1/5})$. Counting the entire rank-four lattice
already gives $O(N^{4/5})$ exceptional inputs. The rational-value
locus in Proposition~\ref{prop:quintic-discard} has dimension one,
which improves this to $O(N^{1/5})$. Both counts are global: the
exceptional set is removed once before any box, prime partition or
curve is selected.

\paragraph{A box-incidence estimate.}
The elementary box-incidence estimate can be stated in a concrete grid.
Fix $C\ge1$, $M>T\ge1$, and grid bases
$c=(c_1,c_2,c_3)\in\Z^3$ with $|c_i|\le CM$. Call a nonzero vector
$v=(v_0,v_1,v_2,v_3)\in\Z^4$ a witness for $c$ if
\[
 |v_0|\le T,\qquad |Mv_i-c_i v_0|\le T\quad(1\le i\le3).
\]
Then $v_0\ne0$, since $v_0=0$ and $M>T$ would force every $v_i=0$.
Moreover $|v_i|\le C|v_0|+T/M\le(C+1)|v_0|$, so, with the maximum
norm, $\|v\|_\infty\asymp_C|v_0|$ and $\|v\|_\infty\ll_C T$.
For fixed $v$, each $c_i$ has $O_C(T/|v_0|)$ possible integer values.
Consequently $v$ witnesses at most
\[
 O_C\bigl((T/\|v\|_\infty)^3\bigr)
\]
grid boxes. These are original coordinate vectors; the residual vector
$(v_0,Mv_1-c_1v_0,Mv_2-c_2v_0,Mv_3-c_3v_0)$ used in
Proposition~\ref{prop:lattice-boxes} is a different vector.

In the density proof,
Proposition~\ref{prop:lattice-boxes} counts boxes with its single prime
approximation grid, and Theorem~\ref{thm:affine-counting} handles adaptive
auxiliary primes separately inside each transformed affine box.

The exact density follows from the general large-prime estimate and
the finite sieve. The Euclidean factorization, level equation and Pell
estimate use the explicit arithmetic of $T^4+1$. The analytic-relation
estimate is conditional on its stated hypotheses.
